\section{Region Labeling}
\subsection{Naive region labeling}
My implementation of the naive algorithm can be seen below.
\begin{lstlisting}[language={futhark}]
def region_label_naive [h] [w] (img: [h][w]u32) : [h][w]i64 =
  let labels_flat = flatten (tabulate_2d h w \i j -> flat_pos w (i, j))
  let (_, res) =
    loop (prev_labels, current_labels) = (replicate (h * w) (-1), labels_flat)
    while prev_labels != current_labels do
      let new_labels =
        map (\fpos ->
               let pos = unflat_pos w fpos
               let color = get pos img
               let neighbours =
                 map (\d ->
                        let npos = move d pos
                        in if in_bounds img npos
                           then (npos, get npos img)
                           else (no_pos, 0))
                     [#n, #w, #e, #s]
               let same_color_labels = map (\(npos, c) -> if c == color && npos != no_pos then current_labels[flat_pos w npos] else -1i64) neighbours
               in reduce i64.max current_labels[fpos] same_color_labels)
            labels_flat
      in (current_labels, new_labels)
  in unflatten (res :> [h * w]i64)
\end{lstlisting}

The implementation maps over all the pixels in the images. It checks which
neighbours have the same color, and the neighbours without the same color are
set to -1, such that they are not chosen in the reduce.

Inside the loop we have a map function, which runs over all pixels in the
image. Inside the map, we do a constant amount of work (although we do have a
map it is always over the neighbours and the same with the reduce). As such the
inside of the loop has work O($N$) and span O($1$). However, figuring out how
many times the loop runs is a little more tricky. 

If the entire image is all the same color, then we will have to propagate the
largest label throughout the entire image. This label would be in the bottom
right corner of the image. At each iteration it moves one pixel in each
direction, so before reaching the top left corner, it would have to do $H + W$
iterations, where H is the height and W is the width. As the loop is
sequential, this gives us a total work of O($O \times (H + W)$) and span of
O($H + W$). Note that $H + W - 1 =< H * W = N$\footnote{the -1 comes as in the
example before we should not count the bottom left corner twice.} as both H and
W are greater or equal to 1.

\subsection{Enlightened region labeling}
Below my code for the algorithm can be seen.
\begin{lstlisting}[language={futhark}]
def norm_edge' (a: i64, b: i64) : (i64, i64) =
  if a <= b then (a, b) else (b, a)

-- | Create normalised edges linking all neighbouring pixels with the same
-- colour.
def mk_edges [h] [w] (img: [h][w]u32) : ?[k].[k]edge =
  let directions = [#n, #w, #s, #e]
  let edges =
    tabulate_2d h w (\r c ->
                       let pos = (r, c)
                       let color = get pos img
                       let edges =
                         map (\d ->
                                let neighbour = move d pos
                                in if in_bounds img neighbour
                                   && color == (get neighbour img)
                                   then norm_edge w (pos, neighbour)
                                   else (no_pos, no_pos))
                             directions
                       in edges)
    |> flatten_3d
    |> filter (\(a, b) -> a != no_pos && b != no_pos)
  in edges

def region_label_enlightened [h] [w] (img: [h][w]u32) =
  -- Step 1: compute edges.
  let edges = #[trace] map (\(a, b) -> (flat_pos w a, flat_pos w b)) (mk_edges img)
  -- Step 2: Initialise DAG.
  let forest = flatten (tabulate_2d h w \i j -> flat_pos w (i, j))
  let (forest', _) =
    loop (forest, edges) while length edges > 0 do
      let possible_edges = filter (\(u, _) -> forest[u] == u) edges
      let sources = map (.0) possible_edges
      let targets = map (.1) possible_edges
      let new_forest = reduce_by_index forest i64.max (-1) sources targets
      let non_inserted_edges = filter (\(a, b) -> new_forest[a] != b) edges
      let new_edges = map (\(a, b) -> norm_edge' (new_forest[a], new_forest[b])) non_inserted_edges
      in (new_forest, new_edges)
  in map (\a' -> loop a = a' while a != forest'[a] do forest'[a]) forest' |> unflatten
\end{lstlisting}

% TODO: Have to talk about the function and tests
